\chapter{Keymaps and bindings}
\label{ch:keymaps}

\begin{aside}
A keymap is a function from key event to action. Users want
to remap; widgets want defaults; chord prefixes need state.
A small system handles all three.
\end{aside}

\section{The map}

\begin{lstlisting}
struct Binding {
    KeyEvent key;
    std::function<void()> action;
    std::string description;
};

class Keymap {
    std::vector<Binding> bindings;
public:
    bool dispatch(KeyEvent ev) {
        for (auto& b : bindings)
            if (matches(b.key, ev)) {
                b.action(); return true;
            }
        return false;
    }
};
\end{lstlisting}

\section{Layered keymaps}

Keymaps stack: widget keymap, then panel keymap, then
global keymap. Dispatch tries each in order; first match
wins.

\begin{lstlisting}
class KeymapStack {
    std::vector<Keymap*> layers;
public:
    void dispatch(KeyEvent ev) {
        for (auto* k : layers)
            if (k->dispatch(ev)) return;
    }
};
\end{lstlisting}

\section{Modes}

A modal app (vim-style) has a different keymap per mode.
Switch via mode signal:

\begin{lstlisting}
auto km = ctx.get(mode) == Mode::Insert
    ? &insert_map
    : &normal_map;
\end{lstlisting}

\section{Chord prefixes}

\code{g d} for ``go to definition'' is a chord. The first
key (\code{g}) sets a prefix state; the next key resolves
against the prefix's submap.

\begin{lstlisting}
auto current = bindings;
on_key([&](KeyEvent ev) {
    auto next = current.lookup(ev);
    if (next.is_submap()) current = next;
    else if (next.is_action()) {
        next.action();
        current = bindings;
    }
});
\end{lstlisting}

A timeout resets the prefix if the second key doesn't
arrive in time.

\section{User remapping}

The keymap is data; load from config:

\begin{lstlisting}
[[bindings]]
key = "Ctrl-S"
action = "save"
\end{lstlisting}

Action names map to functions; the loader looks them up by
name.

\section{Help screen}

A help command lists every binding with its description.
Generated from the keymap itself.

\section{Conflict detection}

Two bindings on the same key: the second wins. \maya{}'s
loader warns on conflicts.

\section{Recap}

\begin{recap}
\begin{itemize}[leftmargin=*]
  \item Keymap = list of (key, action, description).
  \item Layered stack for widget/panel/global.
  \item Modes swap maps.
  \item Chords via prefix state.
  \item User remap via config; help screen lists.
\end{itemize}
\end{recap}

\section*{Workshop}

\begin{enumerate}[leftmargin=*]
  \item Build a keymap for your app's main commands.
  \item Add a chord binding (e.g., \code{Space f} for
    file picker).
  \item Implement a help screen that lists bindings.
\end{enumerate}
